\originalsection{把电容网络还原为节点}\label{tb:cap-networks}
判断串并联要看导线连接的节点, 而不是图上元件的位置. 串联节点不能与第三条导电支路相连, 并联元件则必须有完全相同的两端节点. 原例5.4是$C_1,C_2$先并联, 再与$C_3$串联. 先得$C_{12}=C_1+C_2$, 再得$C_{\rm eq}=(C_1+C_2)C_3/(C_1+C_2+C_3)$. 外加电压$V$时, 串联块的电荷$Q=C_{\rm eq}V$, 并联块电压$Q/(C_1+C_2)$, 两只的电荷按$C_1:C_2$分配, 而$C_3$独自带$Q$.

\begin{example}
\textbf{原5.10.1: 一、二、三只电容组成的三支路.} 
六只相同电容$C$接成三条并联支路: 第一条一只, 第二条两只串联, 第三条三只串联. 求总电容及接电压$V$后各只的电荷、电压.
\end{example}
\begin{solution}
三支路电容是$C,C/2,C/3$, 总量$C_{\rm eq}=11C/6$. 第一支路电荷$CV$, 第二支路各只电荷$CV/2$且各分$V/2$, 第三支路各只电荷$CV/3$且各分$V/3$. 总电源正端供电量是三支路的$CV+CV/2+CV/3$, 不能把六只正极板电荷全部相加, 因为内部浮置节点上的正负电荷来自重排而非都经过电源.
\end{solution}

\originalsection{电场能量的量级与来去}\label{tb:energy-create}
\begin{example}
\textbf{原5.5、5.6: 空气极限与带电球壳.} 
干燥空气击穿场以$3.0\times10^6\,\mathrm{V/m}$估计, 求对应场能密度. 半径$R$带电薄导体球壳的总电荷为$Q$, 求总场能.
\end{example}
\begin{solution}
第一问$u_E=\eps E^2/2\simeq39.8\,\mathrm{J/m^3}$. 击穿场是近似实验量, 会随气压、几何和杂质变化. 第二问壳内场为零, 外部$E=\kc Q/r^2$, 所以
\[U=\int_R^\infty\frac\eps2\frac{\kc^2Q^2}{r^4}4\pi r^2\dd r
=\frac{Q^2}{8\pi\eps R}=\frac{Q^2}{2C},\quad C=4\pi\eps R.\]
比均匀实心球自能小, 因为导体电荷移到表面改变了分布. 当$R\to0$时自能发散, 有限球壳不能直接成为有限自能的经典点电荷.
\end{solution}

原5.6关于创造偶极和消灭场能的动画可由装配功理解. 两个相反电荷从远处分离状态缓慢靠近, 相互作用能$U_{\rm int}=-\kc q^2/r$下降, 场对外部约束做正功; 外力若缓慢移动, 必须取走这部分能量. 但各电荷的自能仍在, 所以不能把负的相互作用能称为负的总场能. 两同号电荷靠近则需要外界输入功. 原件动画的变形图像在这里对应明确的能量差.

\originalsection{极化电荷与两种分层介质}\label{tb:dielectric-boundary}
极化强度$\vct P$是单位体积的偶极矩, 单位C/m$^2$. 一小体积中正负束缚电荷偏移后, 邻近体积的电荷通常抵消, 只在极化不均匀处或表面留下净量. 定量关系是
\[\rho_b=-\nabla\cdot\vct P,\qquad \sigma_b=\vct P\cdot\uv n.\]
这里$\uv n$是介质向外的法向. 均匀极化块内无体束缚电荷, 相对表面仍有异号电荷. 分子取向的永久偶极和外场诱导的偶极都可导致极化, 不要求每种绝缘体的分子原先都有永久偶极.

对跨界扁柱, $D_{n,2}-D_{n,1}=\sigma_f$只计自由面电荷; 静电切向$E$连续. 无自由界面电荷的沿场叠层中法向$D$相同, 而$E_i=D/\varepsilon_i$不同. 横向并排层两端电压相同、同厚度, 场相同而自由电荷不同. 导体、自由电荷和束缚电荷必须分别辨认.

\begin{example}
\textbf{原5.10.2及分层变式: 厚度方向与面积方向的两种填充.} 
面积$A$、板距$d$的电容器, 两种介质相对介电常数$\kappa_1,\kappa_2$. (a)各占厚度$d/2$; (b)各占面积$A/2$. 忽略边缘, 求电容.
\end{example}
\begin{solution}
(a)同一自由电荷$Q$的通量穿过两层, $V=Qd(1/\kappa_1+1/\kappa_2)/(2\eps A)$, 故
\[C_{\rm series}=\frac{2\eps A}{d}\frac{\kappa_1\kappa_2}{\kappa_1+\kappa_2}.\]
(b)两块均跨全板距, 共同电压$V$, $Q_i=\eps\kappa_i(A/2)V/d$, 所以
\[C_{\rm parallel}=\frac{\eps A}{2d}(\kappa_1+\kappa_2).\]
相同材料体积不保证相同电容, 两者在$\kappa_1=\kappa_2$时才必然一致.
\end{solution}

\begin{example}
\textbf{原5.10.3: 同心球间两种介质.} 
内、外电极半径为$a,c$. 介质1填$a<r<b$, 介质2填$b<r<c$, 常数为$\kappa_1,\kappa_2$. 求电容. 这里中间区域是绝缘介质, 并非一整块导体壳的内部材料.
\end{example}
\begin{solution}
内电极电荷$Q$, 故$D_r=Q/(4\pi r^2)$, 两段势降相加:
\[V=\frac Q{4\pi\eps}\left[\frac1{\kappa_1}\left(\frac1a-\frac1b\right)+\frac1{\kappa_2}\left(\frac1b-\frac1c\right)\right].\]
所以
\[C=\frac{4\pi\eps\kappa_1\kappa_2abc}{\kappa_2c(b-a)+\kappa_1a(c-b)}.\]
令两常数相等还原均匀介质球电容, 令$b=a$消去第一层. 原题措辞易让“球电容器”误读为单个实心导体, 此处以电极和介质的实际角色明确条件.
\end{solution}

\originalsection{电容力、弹簧与固定电压的功}\label{tb:cap-force}
\begin{example}
\textbf{原5.10.4: 板间吸力拉长弹簧.} 
平行板面积$A$, 一板接劲度系数$k$的弹簧, 两板各带固定电荷$\pm Q$. 忽略边缘及重力在运动方向的分量. 求静态弹簧伸长量.
\end{example}
\begin{solution}
间距为$d$, 场能$U_E=Q^2d/(2\eps A)$. 若移动板向另一板靠近$x$, 则$d=d_0-x$, 总势为$k x^2/2+Q^2(d_0-x)/(2\eps A)$. 导数为零得
\[x=\frac{Q^2}{2k\eps A}.\]
要求$x<d_0$且板未接触. 同样结果来自$F=p_EA=Q^2/(2\eps A)$. 直接用整张对面的场$Q/(2\eps A)$作用于$Q$也得同一式.

若改接恒压电源, 这是另一个问题: $F=\eps A V^2/(2d^2)$. 电源提供$V\dd Q$, 场能变化只有$\dd(CV^2/2)$; 作机械功时应使用包含电源交换的有效势$U_{\rm eff}=kx^2/2-C(x)V^2/2$. 单独对正的$CV^2/2$求导会把力方向算反. 此变式说明为何“固定电荷”是原解不可省的条件.
\end{solution}

线性无损介质储能密度为$u=\tfrac12\ee\cdot\vct D=\varepsilon E^2/2$, 比真空式多计极化响应. 非线性介质应沿实际建立过程积分$\int\ee\cdot\dd\vct D$; 有滞后和耗散时一条常数介电率曲线不能描述整个循环能量.

\originalsection{原子结合、偶极作用与静电模型的限度}\label{tb:atomic-forces}
原5.8附录的离子和范德瓦耳斯作用, 是把宏观电场连接到物质结合的入口. 异号离子远距离库仑吸引给负势能$-\kc q^2/r$, 短距离不能让两个有限电子云任意重合; 若以经验排斥$A/r^n$($A>0,n>1$)模拟, 则
\[U(r)=\frac A{r^n}-\frac{\kc q^2}r,\quad
r_0^{n-1}=\frac{nA}{\kc q^2},\quad U''(r_0)=\frac{(n-1)\kc q^2}{r_0^3}>0.\]
这里$r_0$是这个相对距离模型的稳定平衡, $A,n$不是由本章经典电学推出的微观常数. 偶极间的固定取向势可写
\[U_{12}=\frac{\kc}{r^3}\big[\vct p_1\cdot\vct p_2-3(\vct p_1\cdot\uv r)(\vct p_2\cdot\uv r)\big],\]
由一个偶极的势对另一小偶极展开得到. 临时涨落、诱导偶极及取向平均能导致范德瓦耳斯吸引, 其微观系数需要更进一步的统计或量子知识; 静止的两个中性点没有自动的库仑单极吸引.

二维或三维电荷阵列、偶极晶格、悬桥和粒子阱演示都需说明可动自由度. 在无电荷区域$\nabla^2V=0$, 若$qV$有非退化三维极小, Hessian矩阵正定将使三个主曲率之和为正, 与拉普拉斯方程矛盾. 对可能退化的严格局部极小, 还要用调和函数的球面平均性质: 对固定测试电荷$q$,$U=qV$也为调和函数, 无源小球中心的$U$等于球面平均$U$. 若球面各点的$U$都比中心更高, 平均值也更高, 同样矛盾, 对负$q$仍成立. 这个平均性质可由各半径球面的高斯零通量导出球面均值对半径不变, 并在半径趋零时等于中心值. 因而仅靠静电场不能稳定悬浮一个自由点电荷于三维局部极小; 固定轨道、支撑杆、时间变场或其他力会改变这个结论的假设. 图中的二维稳定不能据此认定离开平面也稳定.

\begin{example}
\textbf{原5.7: 任意位置的厚介质板.} 平行板面积$A$, 板距$d$, 介质厚$t\le d$, 相对介电常数$\kappa$, 介质覆盖整个面积但不必居中. 求电容及三个极限.
\end{example}
\begin{solution}
无自由界面电荷, 各层$D=Q/A$相同. 真空总厚度$d-t$, 介质厚度$t$, 总势差$V=Q(d-t+t/\kappa)/(\eps A)$. 因而
\[C=\frac{\eps A}{d-t+t/\kappa}.\]
$t\to0$或$\kappa\to1$回到空气电容$\eps A/d$, $t\to d$给完全填充$\kappa\eps A/d$. 介质位置不影响各段总厚度之和, 在平行板、无边缘和线性条件下电容不随其沿法向平移改变.
\end{solution}

拉开两相反电荷的动态过程还含传播与辐射. 移动电荷时, 场的局部能量交换率为$q\vct v\cdot\ee_{\rm other}$, 正值表示场对电荷增加机械能, 负值表示电荷把能量送给场. 电荷停止运动以后, 新的静电构型不能立即在全空间出现; 它以有限速度向外建立, 波前外暂保留旧场. 加速阶段产生的向外辐射可在源已静止后继续离开, 最终留下新静场及远去的辐射能. 因此原5.6“创造/消灭场”的动画须与静态装配功区别, 完整能量账还包括约束力作功和离开的辐射. 第\ref{ch:maxwell}章将给局部能流方程来量化这件事.

原附录的可交互电荷/偶极阵列还规定了额外短程排斥和速度阻尼, 让重叠被抑制、运动逐步耗散到稳定构型. 粒子阱演示另施加与电荷符号无关的向心约束, 悬桥演示固定端点并加重力. 这些模拟的稳态并非单靠库仑力实现自由三维静电悬浮, 所以不能拿一幅稳定阵列图推翻上述调和势结论. 本稿说明各模型的附加力, 不把未实际运行的原交互程序声称为已复现数值实验.
